\section{The Frobenius bound in the integral lifting argument}

% BEGIN F84R1
\paragraph{F84R1. A uniform power of the variable modulo a power of the prime.}
Let $A=O[[u]]$, let $\phi$ fix the coefficients and replace
$u$ by $u^p$, and let $I_\tau$ be ideals indexed cyclically
by $\tau\in\mathbb Z/a$ with $a\ge1$. Write
$E=u^e+p$. Suppose $1\le b$, $h+b\le p$, each
$I_\tau+(p^b)$ contains a power of $u$, and
$E^h I_{\tau+1}\subset A\phi(I_\tau)$.
Then, simultaneously at every $\tau$,
\[
               u^{e(h+b)}\in A\phi(I_\tau)+(p^b).
\]
To prove this, write each series uniquely as
$\sum_{j=0}^{p-1}u^j\phi(f_j)$ by grouping its coefficients
modulo $p$. If a series belongs to $A\phi(I)+(p^b)$ then
each $f_j$ belongs to $I+(p^b)$: for a term
$g\phi(z)$ the $j$th component is $g_j z$, and the operation
commutes with multiplication by $p^b$. Let $s_\tau$ be the
least exponent with $u^{s_\tau}\in I_\tau+(p^b)$.
In $(E-p)^{h+b-1}$ each term is divisible either by $E^h$
or by $p^b$. Multiplying by $u^{s_{\tau+1}}$ therefore puts
$u^{s_{\tau+1}+e(h+b-1)}$ in $A\phi(I_\tau)+(p^b)$.
Take its single nonzero component in the residue-class
decomposition. We get
\[
        p s_\tau\le s_{\tau+1}+e(h+b-1).
\]
At a maximal $s_\tau$ this gives
$(p-1)s_\tau\le e(h+b-1)\le e(p-1)$, hence
$s_\tau\le e$ at every index. Multiplying the preceding
monomial by the remaining nonnegative power of $u$ proves
the assertion. This argument uses no preliminary numerical
bound on the nilpotence exponent of $u$.
% END F84R1

% BEGIN F84R2
\paragraph{F84R2. The first approximation on the actual Frobenius lattice.}
Consider the integral models of Q84T3 with their original
heights $[-\eta,H]$. For $G:M_1\to M_2$ the given
map write $G_\tau$ for its component matrix in integral bases,
$A_{j,\tau}$ for the Frobenius matrices and $L_{j,\tau}=
A_{j,\tau}M_{j,\tau}$ for their column lattices in those
coordinates. The matrices satisfy
$G_\tau A_{1,\tau}=A_{2,\tau}\phi(G_{\tau-1})$,
and the height inequalities are
$E^H M_j\subset L_j$ and $E^\eta L_j\subset M_j$.
For $1\le b$ with $H+\eta+b\le p$ one has
\[
       (u^e)^{H+\eta+b} L_2\ \subset\ G L_1+p^b L_2 .
\]
Set $I_\tau=\operatorname{Ann}_A\operatorname{coker}(G_\tau)$.
The finite-length cokernel of Q84T3 implies that a power of $u$
lies in each $I_\tau$, since $u$ belongs to the maximal ideal.
The annihilator for the matrix $\phi(G_\tau)$ is exactly
$A\phi(I_\tau)$. Indeed apply the component decomposition of
F84R1 entrywise to an equation
$\phi(G_\tau)y=d\phi(x)$ for each integral $x$. The $j$th
component shows $d_j x\in\operatorname{im}G_\tau$.
Conversely, matrices of preimages of the multiples of standard
basis vectors by any $d_j\in I_\tau$ give preimages for
$\phi(d_j)$ after applying $\phi$, proving the other inclusion.

For $d\in I_\tau$ and any integral vector $z$ at the target
of the Frobenius coordinates, $E^\eta A_{2,\tau} z$ is integral.
Choose $m\in M_1$ with $G_\tau m=d E^\eta A_{2,\tau}z$.
Write $E^H m=A_{1,\tau}m'$ with $m'$ integral. Commutation with
Frobenius gives
$A_{2,\tau}\phi(G_{\tau-1})m'=
 A_{2,\tau}dE^{H+\eta}z$. Cancel $A_{2,\tau}$ over the
fraction field. The preceding annihilator identity gives
$E^{H+\eta}I_\tau\subset A\phi(I_{\tau-1})$.
Apply F84R1 with $h=H+\eta$. Multiplication by
$u^{e(H+\eta+b)}$ sends each vector $z$ into
$\phi(G_{\tau-1})M_1+p^b M_2$ in these integral coordinates.
Multiplying by $A_{2,\tau}$ and commuting $G$ back proves the
display on all components. It is a bound on a power of $u$
modulo $p^b$; the subsequent monodromy calculation is needed
to get a bound on a power of $p$.
% END F84R2

% BEGIN F84R3
\paragraph{F84R3. Passing a sharp cokernel bound to all the invariant lattices.}
In the same coordinates suppose that $p^w L_2\subset G L_1$.
Then $p^w M_2\subset G M_1$, and the same exponent works on
$E^{-\eta}M$, $M^\Delta$, $Y=(E^{-\eta}M)^\Delta$ and
$Y'=u^e Y$. Here $\Delta$ is the tame group used for the
invariants. For $z\in M_{2,\tau-1}$,
$A_{2,\tau}\phi(z)\in L_{2,\tau}$. Lift its $p^w$ multiple
to $A_{1,\tau}x$ with $x$ integral. Canceling $A_{2,\tau}$
as above gives $\phi(G_{\tau-1})x=p^w\phi(z)$.
On power series put
$\psi(\sum c_j u^j)=\sum c_{pj}u^j$. It preserves
integrality and satisfies $\psi(\phi(g)d)=g\psi(d)$ by direct
coefficient comparison. Thus $G_{\tau-1}\psi(x)=p^w z$,
giving the claim on $M$ without dividing the exponent by the
Frobenius degree. Scaling lifts by $E^{-\eta}$ gives the
next lattice. For an invariant target, average a lift over
$\Delta$; its order is a unit in $O$, and equivariance shows
that the average has the same image. Multiplication by $u^e$
then handles $Y'$. These implications isolate a single sharp
cokernel estimate on $L$ in the lifting calculation of Q84T4.
% END F84R3
